\section{Related Work}
\label{sec:related}

\subsection{Network Slicing and Resource Allocation}

Network slicing has been studied extensively since its
standardization in 5G, from architecture
frameworks~\cite{rost2017network,afolabi2018slicing} to
orchestration and lifecycle management~\cite{foukas2017networking}.
Resource allocation for slices spans a wide methodological spectrum.
At the static end, operators size slices from engineering rules and
peak-load studies; these approaches are transparent and stable but
cannot track demand dynamics, and empirical studies repeatedly report
substantial over-provisioning as the price of
SLA safety~\cite{sciancalepore2017mobile}. At the dynamic end,
optimization-based approaches formulate slice embedding, scaling, or
admission as integer programs solved on rolling
horizons~\cite{yu2008rethinking}, and learning-based approaches
train reinforcement-learning agents to adjust slice resources from
feedback~\cite{wang2017intelligent,bega2019deepcog}. RL-based slice controllers
can in principle discover policies beyond human design, but they
raise adoption barriers that this paper takes seriously: the policy
is a black box whose failure modes are hard to audit, the training is
sample-hungry on infrastructure where exploration is costly, and the
learned controller is notoriously sensitive to distribution shift.
UCRA deliberately positions itself between the extremes: the forecast
is learned, but the decision layer is a single closed-form transform
that an operator can read, clamp, and reason about, in the spirit of
interpretable control for safety-critical
infrastructure~\cite{rudin2019stop}.

\subsection{Demand Prediction for Network Operations}

Traffic prediction has a long history in networking, from
classical time-series models~\cite{box2015time} to deep sequence
models. Recurrent architectures, and LSTMs in
particular~\cite{hochreiter1997lstm}, remain strong, well-understood
baselines for 15-minute to hourly network demand forecasting, with
convolutional and attention alternatives offering accuracy gains on
dense multivariate settings~\cite{zhang2017deep}. Applications
include cellular traffic forecasting for proactive
scaling~\cite{zhang2019transfer}, capacity planning, and anomaly
detection. Point forecasts, however, answer only part of the
reservation question: committing capacity requires knowing how wrong
the forecast can be. Probabilistic alternatives include Bayesian
neural networks~\cite{gal2016dropout}, deep ensembles, and --- most
directly relevant here --- \emph{quantile regression}, which trains a
model to output several conditional quantiles simultaneously by
minimizing the pinball (quantile) loss~\cite{koenker1978regression}.
Quantile forecasts are standard in electricity markets and load
serving~\cite{hong2016probabilistic} precisely because planners must
commit resources against uncertain demand, and we import that
tradition into slicing. Calibration --- whether the nominal 90\%
interval actually covers ~90\% of outcomes --- is the property that
makes such forecasts decision-grade, and we measure it explicitly
(Section~\ref{sec:results-calibration}).

\subsection{From Forecasts to Actions: Predictive Resource Management}

Converting predictive uncertainty into resource decisions has been
explored in cloud and network contexts. Forecast-driven autoscaling
sizes server fleets from predicted load, with queueing-theoretic
headroom rules~\cite{gandhi2012autoscale};
probabilistic variants size from predicted demand
distributions~\cite{ islam2012provisioning}. In wireless networks,
predictive resource allocation uses traffic forecasts to pre-allocate
PRBs, pre-warm carriers, or schedule sleep states. Closest in spirit
to UCRA, measurement-based admission control sized capacity from
recent traffic statistics with robustness
guarantees~\cite{grossglauser1999adaptive}.
Three gaps remain. First, most of this literature treats the
forecast-to-action mapping as a fixed engineering constant (a safety
factor), losing the connection between the operator's risk appetite
and the frontier position. Second, the risk feedback loop --- using
recent reservation outcomes to modulate the buffer --- is usually
absent. Third, the evaluation rarely closes the loop: metrics are
computed on forecast accuracy (pinball loss, coverage) rather than on
the operational reservation outcomes (violations, utilization, waste)
that the deployment actually cares about. UCRA addresses all three:
$\Phi$ exposes a single risk dial $\kappa$ whose frontier we trace
empirically; the risk indicator $\rhot$ closes the feedback loop
inside the transform; and all reported metrics are closed-loop
reservation outcomes on held-out windows.

\subsection{Drift Detection and Continual Learning}

Distribution shift is the normal condition of production networks:
traffic mix evolves, events recur, application behavior changes.
Drift detection is a mature field --- sequential tests such as
Page's CUSUM and the Page--Hinkley test~\cite{page1954continuous},
ADWIN's adaptive windowing~\cite{bifet2007learning}, and
population-stability statistics are standard tools. On detection,
continual-learning methods update the model without erasing prior
competence; experience replay is the simplest and most robust family,
and reservoir sampling~\cite{vitter1985random} maintains a uniform
sample of history in a fixed-size buffer, which directly mitigates
catastrophic forgetting~\cite{chaudhry2019continual,rolnick2019experience}.
UCRA's Stage~5 combines a dual-trigger monitor (rolling reservation
violation rate OR demand z-score) with reservoir replay and gentle
fine-tuning. What is novel is not the ingredients but the
\emph{contract}: a formally stated information-flow rule for what the
updater is allowed to know, enforced in code and tested, motivated by
the observation that online-learning evaluations routinely --- and
often unknowingly --- grant the learner access to labels that do not
yet exist, a special case of the leakage taxonomy of
Kapoor and Narayanan~\cite{kapoor2023leakage}. Prequential
evaluation~\cite{dawid1984present} formalizes the honest variant:
predict-then-observe in strict temporal order; our protocol is the
multi-step-horizon generalization of that principle for replay-based
updaters.

\subsection{Positioning Against Learned Slice Controllers}

The closest learned-controller line of work is DeepCog
\cite{bega2019deepcog}, which trains a deep network to output slice
capacity demand forecasts that a downstream optimizer converts into
slicing decisions, reporting large gains in resource efficiency on
synthetic and emulated topologies; mobile-traffic-forecast-driven
slicing utilization was likewise analyzed by Sciancalepore et
al.~\cite{sciancalepore2017mobile}. UCRA differs in three deliberate
ways. First, the forecast is \emph{probabilistic} (a full quantile
vector) rather than a point estimate, so the uncertainty the decision
layer must absorb is measured, not assumed. Second, the
decision layer is an explicit equation with one risk parameter rather
than a learned policy or an optimizer over a learned model --- weaker
in raw fitting power, stronger in auditability, and immune to the
distribution-shift brittleness that end-to-end controllers exhibit.
Third, UCRA's evaluation protocol holds the \emph{online-learning
component} to a causal information-flow contract
(Section~\ref{sec:stage5}); learned-controller evaluations, including
ours originally, rarely do, and Section~\ref{sec:leakage-audit}
quantifies what that omission was worth in our own pipeline.

\subsection{Uncertainty-Aware Provisioning in Adjacent Domains}

The pattern that UCRA instantiates --- forecast a distribution, then
commit a resource against a chosen quantile with an explicit risk
dial --- is mature in adjacent domains, and the borrowings are
deliberate. In electricity systems, probabilistic load forecasting
has fed unit commitment and reserve sizing for a decade: system
operators procure spinning reserve against forecast quantiles, and
the literature treats calibration as a procurement-relevant property
exactly as we treat coverage here~\cite{hong2016probabilistic}. In
cloud capacity management, AutoScale provisions server tiers from
predicted load with robustness margins and switching-cost
awareness~\cite{gandhi2012autoscale} --- the direct ancestor of our
Stage~4 smoothing and release hysteresis --- and empirical
provisioning models forecast workload to adapt resource
allocations~\cite{islam2012provisioning}. Measurement-based admission
control established the principle of committing capacity from
recent traffic statistics with robustness
guarantees~\cite{grossglauser1999adaptive}. What the slicing domain
has lacked is the assembly: these ingredients (calibrated
conditional quantiles, a priced buffer, feedback modulation,
churn-bounded actuation, and honest online adaptation) exist
separately, but no slicing resource manager composes them into one
auditable loop, evaluated closed-loop on live network data. That
composition is UCRA's contribution, and the adjacent-domain
literature provides the confidence that each ingredient is sound.

\subsection{Datasets and Reproducibility in 5G Research}

Progress in this area depends on public, realistic datasets. The CTTC
B5G slicing dataset~\cite{farreras2024cttc} provides steady-state
slicing simulation snapshots with per-slice reservations, offered
traffic, and QoS outcomes; the live RAN performance-counter
dataset of Lehoczky et al.~\cite{lehoczky2026ran} provides three
weeks of 15-minute sector-level 5G/4G/2G counters from an operational
network. Both are on Zenodo with stable identifiers and are used
here. Reproducibility has become a pressing concern for ML-for-networks
research~\cite{pineau2021improving}; UCRA releases its full pipeline
--- ingestion, training, closed-loop evaluation, causality unit
tests, and the three-seed verification driver --- with pinned
configuration files so that every number in Section~\ref{sec:results}
can be regenerated from a clean clone.
